\section{Regulación: la Ley de IA de la Unión Europea}

\subsection{Panorama general de la EU AI Act}

El expositor presentó en detalle la Ley de Inteligencia Artificial de la Unión Europea (EU AI Act), el primer marco jurídico integral de regulación de la IA a nivel mundial.

\begin{importantbox}{Elementos centrales de la EU AI Act}
\begin{itemize}
    \item \textbf{Entrada en vigor}: aprobada formalmente en agosto de 2024, con implementación por etapas que otorga a las empresas un periodo de transición para el cumplimiento
    \item \textbf{Principio central}: regulación escalonada basada en el \textbf{Risk Level} (nivel de riesgo)
    \item \textbf{Lógica regulatoria}: cuanto más alto es el riesgo de un proyecto de IA, más estrictas son las restricciones que enfrenta
    \item \textbf{Nivel máximo}: algunas aplicaciones quedan directamente \textbf{prohibidas} (Banned)
\end{itemize}
\end{importantbox}

\subsection{Niveles de riesgo y aplicaciones prohibidas}

La EU AI Act clasifica las aplicaciones de IA en varios niveles de riesgo y aplica distintos grados de exigencia regulatoria a cada uno. Entre ellos, las siguientes aplicaciones quedan explícitamente \textbf{prohibidas}:

\begin{table}[H]
\centering
\caption{Tipos de aplicaciones de IA prohibidas en la EU AI Act}
\begin{tabular}{@{}ll@{}}
\toprule
\textbf{Aplicación prohibida} & \textbf{Descripción} \\
\midrule
Biometría en espacios públicos & Reconocimiento facial en tiempo real para vigilancia pública \\
Sistemas de puntuación social & Calificación crediticia social basada en características de individuos o grupos \\
IA manipuladora & Explota vulnerabilidades humanas para manipulación subliminal \\
Vigilancia predictiva & Predicción delictiva basada únicamente en perfiles \\
\bottomrule
\end{tabular}
\end{table}

El expositor mencionó en particular los \textbf{sistemas de puntuación social} (Social Scoring Systems): sistemas que asignan una calificación numérica a individuos o grupos, calificación que puede usarse en decisiones de alto impacto como la aprobación de crédito o la selección de personal. Señaló que en Estados Unidos ya existen sistemas similares, por ejemplo algoritmos usados para predecir la tasa de reincidencia (Recidivism), que predicen la probabilidad de reincidencia de una persona a partir de datos de sus características. Bajo el marco de la EU AI Act, este tipo de sistemas se consideraría ilegal.

\begin{warningbox}{El problema del sesgo en los sistemas de predicción de reincidencia}
Los sistemas estadounidenses de predicción de reincidencia (como COMPAS) han sido ampliamente criticados por presentar sesgo racial: sobrestiman de manera sistemática la probabilidad de reincidencia de las personas afroestadounidenses y, al mismo tiempo, subestiman la de las personas blancas. Es un caso representativo de cómo un ``algoritmo aparentemente objetivo'' termina amplificando en la práctica los sesgos sociales existentes. La EU AI Act clasifica explícitamente este tipo de sistemas como prohibidos.
\end{warningbox}

\subsection{El requisito de interpretabilidad y el dilema técnico}

La EU AI Act contiene una disposición que, a juicio del expositor, resulta problemática: el requisito de \textbf{Explainability} (interpretabilidad). La ley exige que los sistemas de IA sean capaces de explicar su proceso de decisión: por ejemplo, ¿por qué un automóvil autónomo se detiene? ¿Por qué da vuelta?

El expositor analizó la dificultad de este requisito desde dos niveles:

\begin{enumerate}
    \item \textbf{La explicación en el nivel matemático sí es posible}: podemos describir con precisión cada paso del cómputo de una red neuronal: los valores de activación multiplicados por la matriz de pesos, la aplicación de la función de activación, la comparación con el sesgo. Pero este claramente no es el nivel de explicación que exige la regulación.
    \item \textbf{Una explicación comprensible para los seres humanos es (en teoría) imposible}: los sistemas de aprendizaje profundo son sistemas complejos (Complex Systems); poseen emergencia (Emergence), autoorganización (Self-organization) y adaptabilidad (Adaptivity). Para este tipo de sistemas, dar la razón de una decisión en un nivel de abstracción comprensible para los seres humanos puede ser \textbf{imposible en principio}.
\end{enumerate}

\begin{importantbox}{La paradoja entre los sistemas complejos y la interpretabilidad}
Los sistemas de aprendizaje profundo son sistemas adaptativos complejos (Complex Adaptive Systems). Los científicos que estudian sistemas complejos saben que el comportamiento de este tipo de sistemas emerge de la interacción de una gran cantidad de componentes simples y no puede reducirse sin más al comportamiento de un solo componente. Por eso, exigir ``explicar por qué la IA tomó esta decisión'' puede equivaler, en esencia, a exigir ``explicar por qué el clima es como es'': podemos ejecutar simulaciones, pero no podemos dar una narrativa causal concisa.
\end{importantbox}

\subsection{Resumen de la sección}

La EU AI Act representa un hito importante en la regulación de la IA a nivel mundial. Adopta una estrategia de regulación escalonada por nivel de riesgo, prohíbe las aplicaciones de más alto riesgo (vigilancia biométrica, puntuación social) y plantea requisitos de cumplimiento estrictos para las aplicaciones de alto riesgo. Sin embargo, el requisito de interpretabilidad enfrenta un desafío técnico fundamental: el comportamiento de los sistemas complejos podría no poder explicarse en un nivel comprensible para los seres humanos.

